\section{Verification Strategy}

\subsection{Veil Proof Structure}

The core proof strategy is to model-check small instances early to expose
specification mistakes, and then use Veil's invariant-checking support to prove
the main safety results. The intended proof shape is:

\begin{enumerate}[leftmargin=1.5em]
\item define the abstract state and transitions;
\item formulate book well-formedness and accounting invariants;
\item use bounded checking to find missing side conditions;
\item strengthen invariants until the main safety properties become inductive.
\end{enumerate}

\subsection{Lean Helper Lemmas}

Several obligations are more naturally handled in Lean than directly inside
Veil. These include:

\begin{itemize}[leftmargin=1.5em]
\item preservation of sortedness under insertion and deletion;
\item simple facts about quantities and residual quantities;
\item arithmetic lemmas needed for conservation of cash and assets.
\end{itemize}

\subsection{Optional Velvet Component}

If the main Veil proof stabilises early enough, a small Velvet implementation
of the matching loop can serve as an executable reference. The point of this
extension is not performance, but a proof-aligned executable artifact whose
postcondition links it back to the Veil transition relation.
